\section{Smart Contract Architecture and Settlement}

\subsection{Non-Custodial Escrow Vault: CypherRollVault.sol}
Central to the trustless operation of CypherRoll is the \texttt{CypherRollVault.sol} smart contract, implemented in Solidity 0.8.24. Unlike centralized platforms where customer deposits are commingled with operating accounts, \texttt{CypherRollVault} enforces non-custodial segregation of player assets on-chain.

\subsection{EIP-712 Structured Data Cryptographic Authorization}
Traditional decentralized gaming architectures force users to submit an on-chain transaction for every single game step, creating unbearable latency and gas consumption. CypherRoll solves this by maintaining fast off-chain ledger balances and settling on-chain withdrawals via \textbf{EIP-712} (Typed Structured Data Hashing and Signing).

\paragraph{EIP-712 Domain Separator}
To prevent cross-chain replay attacks and front-running across different contracts, the vault constructs a deterministic domain separator:
\begin{align*}
\mathcal{D} = \text{keccak256}(\text{abi.encode}( & \text{keccak256}(\text{"EIP712Domain(string name,string version,uint256 chainId,address verifyingContract)"}), \\
& \text{keccak256}(\text{bytes("CypherRollVault")}), \\
& \text{keccak256}(\text{bytes("1")}), \\
& \text{block.chainid}, \\
& \text{address(this)}))
\end{align*}

\paragraph{Withdrawal Struct Typehash}
The struct definition for withdrawals is defined as:
$$\text{WITHDRAWAL\_TYPEHASH} = \text{keccak256}(\text{"Withdrawal(address player,address token,uint256 amount,uint256 nonce,uint256 deadline)"})$$

\subsection{Smart Contract Source Code}
Below is the core implementation of \texttt{CypherRollVault.sol}:

\begin{lstlisting}[language=Solidity,caption={CypherRollVault.sol Smart Contract Implementation}]
// SPDX-License-Identifier: MIT
pragma solidity ^0.8.24;

interface IERC20 {
    function transfer(address recipient, uint256 amount) external returns (bool);
    function transferFrom(address sender, address recipient, uint256 amount) external returns (bool);
    function balanceOf(address account) external view returns (uint256);
}

contract CypherRollVault {
    address public immutable owner;
    address public operatorSigner;
    bool public paused;

    bytes32 public immutable DOMAIN_SEPARATOR;
    bytes32 public constant WITHDRAWAL_TYPEHASH = keccak256(
        "Withdrawal(address player,address token,uint256 amount,uint256 nonce,uint256 deadline)"
    );

    mapping(address => uint256) public userNonces;
    mapping(address => uint256) public tokenVaultReserves;

    event Deposited(address indexed player, address indexed token, uint256 amount, uint256 timestamp);
    event Withdrawn(address indexed player, address indexed token, uint256 amount, uint256 nonce, uint256 timestamp);

    modifier onlyOwner() {
        require(msg.sender == owner, "Only owner");
        _;
    }

    modifier whenNotPaused() {
        require(!paused, "Vault is paused");
        _;
    }

    constructor(address _operatorSigner) {
        require(_operatorSigner != address(0), "Invalid signer");
        owner = msg.sender;
        operatorSigner = _operatorSigner;

        DOMAIN_SEPARATOR = keccak256(
            abi.encode(
                keccak256("EIP712Domain(string name,string version,uint256 chainId,address verifyingContract)"),
                keccak256(bytes("CypherRollVault")),
                keccak256(bytes("1")),
                block.chainid,
                address(this)
            )
        );
    }

    function depositETH() external payable whenNotPaused {
        require(msg.value > 0, "Zero deposit");
        tokenVaultReserves[address(0)] += msg.value;
        emit Deposited(msg.sender, address(0), msg.value, block.timestamp);
    }

    function depositERC20(address token, uint256 amount) external whenNotPaused {
        require(token != address(0), "Invalid token");
        require(amount > 0, "Zero amount");
        bool success = IERC20(token).transferFrom(msg.sender, address(this), amount);
        require(success, "Transfer failed");
        tokenVaultReserves[token] += amount;
        emit Deposited(msg.sender, token, amount, block.timestamp);
    }

    function withdraw(
        address token,
        uint256 amount,
        uint256 nonce,
        uint256 deadline,
        bytes calldata signature
    ) external whenNotPaused {
        require(block.timestamp <= deadline, "Signature expired");
        require(nonce == userNonces[msg.sender] + 1, "Invalid nonce sequence");
        require(tokenVaultReserves[token] >= amount, "Insufficient vault reserves");

        bytes32 structHash = keccak256(
            abi.encode(WITHDRAWAL_TYPEHASH, msg.sender, token, amount, nonce, deadline)
        );
        bytes32 digest = keccak256(abi.encodePacked("\x19\x01", DOMAIN_SEPARATOR, structHash));

        address recovered = recoverSigner(digest, signature);
        require(recovered == operatorSigner, "Invalid operator signature");

        userNonces[msg.sender] = nonce;
        tokenVaultReserves[token] -= amount;

        if (token == address(0)) {
            (bool sent, ) = payable(msg.sender).call{value: amount}("");
            require(sent, "ETH transfer failed");
        } else {
            bool sent = IERC20(token).transfer(msg.sender, amount);
            require(sent, "Token transfer failed");
        }

        emit Withdrawn(msg.sender, token, amount, nonce, block.timestamp);
    }

    function recoverSigner(bytes32 digest, bytes memory sig) internal pure returns (address) {
        require(sig.length == 65, "Invalid signature length");
        bytes32 r; bytes32 s; uint8 v;
        assembly {
            r := mload(add(sig, 32))
            s := mload(add(sig, 64))
            v := byte(0, mload(add(sig, 96)))
        }
        return ecrecover(digest, v, r, s);
    }
}
\end{lstlisting}

\subsection{Security Invariants Enforced by the Vault}
\begin{enumerate}[leftmargin=1.5em]
    \item \textbf{Monotonic Nonce Sequence:} Withdrawals strictly require $\text{nonce} = \text{userNonces}[\text{msg.sender}] + 1$. This eliminates transaction replay attacks and guarantees order integrity.
    \item \textbf{Deadline Expiration:} Signatures expire after 3600 seconds, preventing stale or compromised vouchers from circulating indefinitely.
    \item \textbf{Checks-Effects-Interactions Pattern:} The contract updates reserves and nonces \textit{before} performing external calls, neutralizing reentrancy attacks.
\end{enumerate}
